%============================================================
% CHƯƠNG 2: ĐẠO LÀM BẠN
%============================================================
\chapter{Đạo Làm Bạn (The Way of Friendship)}

\begin{center}
  \textit{\color{truyenblue}``A friend is someone who knows all about you and still loves you.''}\\
  \textit{(Bạn bè là người biết tất cả về bạn mà vẫn yêu thương bạn.)}\\[4pt]
  \small\color{truyengold}--- Elbert Hubbard ---
\end{center}
\ngancach

\section{Chiếc Ghế Trống}
\begin{truyen}{Chiếc Ghế Trống}{Trung Thành}
\chuhoa{K}{ }hi người bạn thân nhất bị tai nạn và phải ngồi xe lăn, cả nhóm bạn cũ dần xa lánh.\\
\textit{(When the closest friend had an accident and had to use a wheelchair, the old group of friends gradually drifted away.)}

Chỉ có một người vẫn đến mỗi sáng thứ Bảy, đẩy xe lăn ra công viên, ngồi bên cạnh uống cà phê và kể chuyện tầm phào như cũ.\\
\textit{(Only one still came every Saturday morning, pushing the wheelchair to the park, sitting beside him, drinking coffee and chatting about trivial things just like before.)}

Người ngồi xe lăn hỏi: ``Sao mày vẫn đến?''\\
\textit{(The one in the wheelchair asked: ``Why do you still come?'')}

``\textbf{Vì tao vẫn là bạn mày. Tai nạn không thay đổi điều đó}.''\\
\textit{(``\textbf{Because I am still your friend. The accident does not change that}.'')}
\end{truyen}
\begin{baihoc}
  Tình bạn thật sự được thử thách không phải lúc vui vẻ mà lúc khó khăn. Người ở lại khi hoàn cảnh xấu đi mới thực sự xứng đáng được gọi là bạn thân.
\end{baihoc}

\section{Sự Thật Khó Nghe}
\begin{truyen}{Sự Thật Khó Nghe}{Trung Thực}
\chuhoa{B}{ }ạn thân nhìn thấy người bạn chuẩn bị ký hợp đồng kinh doanh với kẻ lừa đảo.\\
\textit{(The close friend saw his friend about to sign a business contract with a fraud.)}

Anh biết nói ra sẽ bị nghi ngờ, thậm chí mất bạn – nhưng anh vẫn nói.\\
\textit{(He knew speaking up might cause suspicion and even lose the friendship – but he spoke up anyway.)}

Người bạn nổi giận: ``Mày ghen tị với tao à?''\\
\textit{(The friend got angry: ``Are you jealous of me?'')}

Sáu tháng sau, khi vụ lừa đảo bị vỡ lở, người bạn tìm lại anh và nói: ``\textbf{Mày đã cứu tao. Xin lỗi vì đã không tin}.''\\
\textit{(Six months later, when the fraud was exposed, the friend sought him out and said: ``\textbf{You saved me. I am sorry for not believing you}.'')}
\end{truyen}
\begin{baihoc}
  Người bạn dám nói sự thật dù biết có thể bị ghét là người bạn quý giá nhất. Lời nói thật đắng trong khoảnh khắc nhưng ngọt ngào về lâu dài.
\end{baihoc}

\section{Mượn Tiền}
\begin{truyen}{Mượn Tiền}{Nguyên Tắc}
\chuhoa{B}{ }ạn thân xin mượn một khoản tiền lớn để đầu tư và hứa trả sau ba tháng.\\
\textit{(The close friend asked to borrow a large sum of money to invest and promised to repay it in three months.)}

Người cho mượn biết trong lòng rằng \textbf{nếu mất tiền thì cũng mất bạn}.\\
\textit{(The lender knew in his heart that \textbf{if the money was lost, the friendship would be lost too}.)}

Nhưng anh vẫn cho mượn – không phải vì tiền ít, mà vì anh \textbf{tin vào con người bạn mình, không phải vào khoản đầu tư}.\\
\textit{(But he still lent it – not because the amount was small, but because he \textbf{trusted his friend as a person, not the investment}.)}

Ba tháng sau bạn trả đủ, kèm theo câu: ``Cảm ơn mày đã tin tao.'' – ``\textbf{Tao không cho mày mượn tiền. Tao đầu tư vào tao làm bạn mày}.''\\
\textit{(Three months later the friend repaid in full, adding: ``Thank you for trusting me.'' – ``\textbf{I did not lend you money. I invested in being your friend}.'')}
\end{truyen}
\begin{baihoc}
  Tình bạn và tiền bạc là hai thứ rất dễ làm hỏng nhau. Cách xử lý khéo léo dựa trên lòng tin và sự minh bạch mới là nền tảng để giữ được cả hai.
\end{baihoc}

\section{Cái Ôm Không Nói Gì}
\begin{truyen}{Cái Ôm Không Nói Gì}{Đồng Cảm}
\chuhoa{K}{ }hi người bạn gọi điện lúc nửa đêm, giọng thất thần báo rằng vừa ly hôn xong, anh không biết phải nói gì.\\
\textit{(When his friend called at midnight, speaking distractedly that he had just gone through a divorce, he did not know what to say.)}

Anh hỏi: ``Mày ở đâu?'' – ``Trước cổng nhà mày.''\\
\textit{(He asked: ``Where are you?'' – ``In front of your gate.'')}

Anh mở cửa, không nói một lời, \textbf{chỉ ôm bạn đứng yên dưới mưa một lúc lâu}.\\
\textit{(He opened the door, said not a word, and \textbf{simply held his friend standing quietly in the rain for a long time}.)}

Sau này bạn anh kể: ``\textbf{Đêm đó tao không cần ai nói gì. Tao chỉ cần mày ở đó}.''\\
\textit{(Later his friend told him: ``\textbf{That night I needed no one to say anything. I only needed you to be there}.'')}
\end{truyen}
\begin{baihoc}
  Đôi khi sự hiện diện còn mạnh hơn mọi lời an ủi. Người bạn giỏi không cần có câu trả lời hay – họ chỉ cần đủ dũng cảm để không bỏ đi.
\end{baihoc}

\section{Lá Thư Không Gửi}
\begin{truyen}{Lá Thư Không Gửi}{Tha Thứ}
\chuhoa{S}{ }au một cuộc cãi vã lớn, hai người bạn thân không nói chuyện nhau suốt ba năm.\\
\textit{(After a major quarrel, two close friends did not speak to each other for three years.)}

Một người đã viết rất nhiều lá thư xin lỗi nhưng không bao giờ gửi đi.\\
\textit{(One of them wrote many letters of apology but never sent any of them.)}

Khi nghe tin người kia bị bệnh nặng, anh vội vàng đến bệnh viện, run run đặt xấp thư lên bàn.\\
\textit{(When he heard the other was gravely ill, he rushed to the hospital and tremblingly placed the stack of letters on the bedside table.)}

Người bạn bệnh nhìn đống thư, mỉm cười: ``\textbf{Tao biết mày hối hận. Tao cũng vậy. Chúng ta đã lãng phí ba năm rồi}.''\\
\textit{(The sick friend looked at the pile of letters and smiled: ``\textbf{I knew you regretted it. So did I. We have already wasted three years}.'')}
\end{truyen}
\begin{baihoc}
  Tha thứ sớm một ngày là thêm một ngày được sống trong tình bạn. Sự kiêu ngạo có thể ăn mòn tháng năm tốt đẹp nhất của cuộc đời.
\end{baihoc}

\section{Người Giữ Bí Mật}
\begin{truyen}{Người Giữ Bí Mật}{Trung Tín}
\chuhoa{C}{ }ô đã kể cho bạn thân nghe mọi bí mật đau lòng nhất của đời mình.\\
\textit{(She had told her closest friend every most painful secret of her life.)}

Mười năm sau, khi hai người không còn liên lạc, cô gặp lại nhau tình cờ trong một buổi tiệc.\\
\textit{(Ten years later, after they had lost contact, they met again by chance at a party.)}

Cô lo lắng không biết những bí mật năm xưa có bị đem ra bàn tán.\\
\textit{(She worried whether the old secrets had been gossiped about.)}

Người bạn cũ bước đến, siết tay và chỉ nói: ``\textbf{Tao rất vui vì mày đang sống tốt}.'' – Không một lời nhắc đến quá khứ.\\
\textit{(The old friend stepped forward, squeezed her hand and only said: ``\textbf{I am so glad you are doing well}.'' – Not a single word about the past.)}
\end{truyen}
\begin{baihoc}
  Người xứng đáng được tin tưởng là người biết im lặng kể cả khi tình bạn đã tan vỡ. Đó là phẩm giá cao nhất trong đạo làm bạn.
\end{baihoc}

\section{Bộ Quần Áo Mượn}
\begin{truyen}{Bộ Quần Áo Mượn}{Hào Phóng}
\chuhoa{N}{ }gày đi phỏng vấn quan trọng, bạn không có bộ quần áo tươm tất.\\
\textit{(On the day of an important interview, the friend had no decent outfit to wear.)}

Người bạn thân đang ở tỉnh khác nhận tin nhắn lúc nửa đêm, lập tức gọi điện: ``Size mày bao nhiêu? Tao gửi xe khách sớm mai.''\\
\textit{(The close friend in another province received the message at midnight and immediately called: ``What is your size? I will send it by bus first thing tomorrow morning.'')}

Bộ quần áo đến kịp giờ. Cuộc phỏng vấn thành công.\\
\textit{(The outfit arrived just in time. The interview was a success.)}

Mười năm sau, người kia có việc lớn, anh bay về ngay trong ngày chỉ để đứng bên cạnh.\\
\textit{(Ten years later, when that friend had something important, he flew back the same day just to stand by his side.)}

Đó không phải trả nợ. Đó là \textbf{tình bạn nối tiếp tình bạn}.\\
\textit{(That was not repaying a debt. That was \textbf{friendship continuing friendship}.)}
\end{truyen}
\begin{baihoc}
  Tình bạn đẹp không phải là cân bằng sổ sách. Nó là vòng tuần hoàn của lòng hào phóng, khi mỗi người đến lượt mình đều sẵn sàng cho đi mà không tính toán.
\end{baihoc}

\section{Người Duy Nhất Vỗ Tay}
\begin{truyen}{Người Duy Nhất Vỗ Tay}{Ủng Hộ}
\chuhoa{T}{ }rong buổi biểu diễn đầu tiên của cô, khán phòng im lặng vì màn trình diễn còn nhiều vấp váp.\\
\textit{(At her very first performance, the hall was silent because the show had many stumbles.)}

Hầu hết khán giả nhìn nhau ái ngại. Duy chỉ có một người \textbf{đứng dậy vỗ tay một mình}.\\
\textit{(Most of the audience exchanged awkward glances. Only one person \textbf{stood up and clapped alone}.)}

Là người bạn thân của cô. Anh vỗ đến khi cả phòng cũng vỗ theo.\\
\textit{(It was her close friend. He clapped until the whole room joined in.)}

``Sao mày làm vậy? Tao diễn dở mà.'' – ``\textbf{Mày diễn lần đầu. Tao vỗ cho lần thứ hai của mày}.''\\
\textit{(``Why did you do that? I performed poorly.'' – ``\textbf{It was your first time. I clapped for your second time}.'')}
\end{truyen}
\begin{baihoc}
  Người bạn thực sự không cổ vũ khi ta đã thành công – họ cổ vũ khi ta chưa xứng đáng được cổ vũ. Đó là sức mạnh kỳ diệu của tình bạn chân thành.
\end{baihoc}

\section{Chuyến Đi Không Có Ảnh}
\begin{truyen}{Chuyến Đi Không Có Ảnh}{Hiện Diện}
\chuhoa{H}{ }ai người bạn đi du lịch cùng nhau. Một người liên tục chụp ảnh và đăng mạng xã hội.\\
\textit{(Two friends went on a trip together. One kept taking photos and posting on social media.)}

Người còn lại không chụp một tấm nào. Chỉ ngồi ngắm cảnh, uống trà.\\
\textit{(The other did not take a single photo. Just sat gazing at the scenery, drinking tea.)}

``Mày không chụp ảnh làm kỷ niệm à?''\\
\textit{(``Are you not taking photos as souvenirs?'')}

``\textbf{Tao đang lưu chuyến đi này vào một nơi không ai hack được} – bộ nhớ của tao. Và tao đang lưu cả khoảnh khắc ngồi cạnh mày nữa}.''\\
\textit{(``\textbf{I am saving this trip somewhere no one can hack} – my memory. And I am saving the moment of sitting beside you as well}.'')}
\end{truyen}
\begin{baihoc}
  Hiện diện trọn vẹn là món quà đẹp nhất ta tặng cho bạn bè. Những khoảnh khắc thực sự sống không nằm trên mạng xã hội mà nằm trong trái tim những người đã ở đó.
\end{baihoc}

\section{Tin Nhắn Hỏi Thăm}
\begin{truyen}{Tin Nhắn Hỏi Thăm}{Quan Tâm}
\chuhoa{M}{ }ỗi năm một lần vào đúng ngày sinh nhật, anh nhận được tin nhắn từ người bạn học cũ chưa gặp lại hai mươi năm.\\
\textit{(Once a year on his exact birthday, he received a text from an old classmate he had not seen in twenty years.)}

Chỉ một câu: ``Chúc mừng sinh nhật. Mày ổn không?''\\
\textit{(Just one line: ``Happy birthday. Are you doing okay?'')}

Anh luôn trả lời và họ nhắn qua lại một lúc rồi lại im lặng cả năm.\\
\textit{(He always replied and they exchanged a few messages before falling silent again for a year.)}

Một năm anh bệnh nặng, anh nhận tin nhắn đó và \textbf{lần đầu tiên gọi điện thay vì nhắn lại}.\\
\textit{(One year when he was gravely ill, he received that message and \textbf{for the first time called instead of texting back}.)}

Người bạn cũ lái xe ba tiếng đến ngay hôm đó. \textbf{Thì ra hai mươi năm chưa bao giờ xa}.\\
\textit{(The old friend drove three hours and arrived that very day. \textbf{It turned out twenty years had never really been a distance at all}.)}
\end{truyen}
\begin{baihoc}
  Tình bạn không cần sự tiếp xúc thường xuyên để duy trì. Chỉ cần mỗi người đều giữ trong tim một chỗ cho nhau, khoảng cách và thời gian trở nên vô nghĩa.
\end{baihoc}
